\section{Problem Formulation}

\subsection*{Provisional research question}
Can a management policy select geometric protection during online packing of
BED-BPP \emph{nominal} orders under \emph{synthetic} dimensional uncertainty
(not BED-BPP field noise) and retain more nominal placed volume than
fixed-margin strategies under a common geometric-failure risk criterion?
The G2 v1 freeze used deterministic known growth and is closed for RL under
guarantee. A v2 i.i.d.\ draft separates immediate geometric risk from layout
consequences; a toy enumeration shows a specified local rule can be suboptimal,
but does not justify an RL pilot.

\subsection*{Active contract}
Original sequence; one container equal to the order's own target; online
$p=s=1$; explicit orientations; containment and non-overlap only. Stop when the
chosen protection admits no candidate, or when a realisation violates geometry.
Useful business volume uses \emph{nominal} dimensions of successfully placed
items; artificial protection volume is not merchandise. BED-BPP supplies
nominal sizes, sequence, and target; realisation errors are synthetic and are
not BED-BPP field measurements.

\subsection*{Inactive in the initial scope}
Placement-position noise, box deformation, pile stability, accessibility,
robot trajectories, and recovery after a collision. No latent episode noise
regime is introduced merely to motivate learning.

\subsection*{Information}
\emph{Before} placing item $t$: nominal size of $t$, revealed geometry of
successfully placed items, container size, the public error law, and candidates
computed only from available information. The agent does \emph{not} observe the
current realisation, future errors, the item suffix, or outcomes of unevaluated
alternative protections. Under item-wise independence, history does not identify
the next error.
\emph{After} a successful placement: the simulator reveals the exact realised
dimensions and updates occupancy accordingly. Exact revelation is a simulator
idealisation, not a demonstrated robot sensing claim.

\subsection*{Candidate setting B and reference A}
\textbf{(B)~Candidate (v2 draft).} Realised sizes follow a public generative
law \(R=N(1+\varepsilon)\) (mm; positivity enforced) with synthetic sensitivity
parameters. Default B1: known i.i.d.\ law across items. Axis dependence within
an item is declared at freeze time. Orientation permutes the same realised
vector. The G2 v1 freeze (\(R=N(1+\alpha)\) with public \(\alpha\)) remains
historical harness evidence only.

\textbf{(A)~Conservative reference.} Known support bounds imply a directly
computable protective margin. Under full guarantee, choosing the minimal
sufficient margin is analytically trivial and is not an RL task.

\subsection*{Action and transition}
The management action selects protection $m$ from a small menu. A
\emph{fixed} placement chooser, identical across methods, picks among
candidates whose envelope is
\codepath{nominal_oriented}+\codepath{margin_oriented} with a shared
front-left-bottom corner (no centred protection). Margin changes the candidate
set and thus the pose; geometric failure is evaluated on the \emph{realised}
AABB at that pose.

Three distinct events are used: \codepath{no_candidate} ends with partial
nominal volume; \codepath{envelope_exceeded} means the realisation outgrows the
reserved envelope and is \emph{not} synonymous with geometric failure---if the
realised box neither overlaps revealed boxes nor exits the container, the
episode continues and occupancy is updated with the \emph{realised} geometry;
\codepath{geometric_failure} ends the episode without recovery. Actions are never
revised retrospectively.

\subsection*{Objective and risk}
Primary score \(J_B\): nominal volume over container volume if the episode has
no \codepath{geometric_failure}, else zero. Immediate risk is
\(\Pr(\mathrm{geometric\_failure}\mid\mathrm{pose},G^{\mathrm{rev}},P_\varepsilon)\),
not candidate existence. Coefficients and a PPO reward are not fixed here.
